% ============================================================
%  Rasmus Nielsen, Grundideernes Logik i kort Begreb.
%  Kjøbenhavn: J. H. Schubothes Boghandel, 1870.
%  XIV + 89 pp.
%
%  TRANSCRIPTION (Danish) — from the Google Books scan.
%  Scan: ~/bibliotek/Nielsen, Rasmus/1870-kort-begreb.pdf (118 PDF pp.,
%        Google Books, British Library copy, digitised 2013; GB id
%        r-JZAAAAcAAJ. An identical earlier download sits beside it as
%        Grundideernes_Logik_i_Kort_Begreb.pdf.)
%
%  PAGE MAP (verified against the PDF text layer, every body numeral read
%  in sequence, 2026-08-09; endpoints to be eye-confirmed in batches):
%        printed  1–89   =  PDF + 20   (PDF 21–109), uniform, no gaps
%        roman  I–XIV    =  PDF + 6    (PDF  7–20)
%        Title leaf I (PDF 7; further title-ish leaves PDF 3, 5); BM-stamp
%        verso II; Forord III; blank IV; Indledning V–XIII (folio suppressed
%        on V); Indhold begins on the lower half of fol. XIII and continues
%        on XIV. Folio suppressed on printed p. 1. Eye-verified: PDF 3 = PDF 5
%        is an ornamental wrapper leaf (imprint without year); PDF 7 is the
%        title page proper; back leaves PDF 110–118 are wrapper-style leaves
%        bearing only »1870.« between rules — no advertisements. The numeral
%        of printed p. 58 (PDF 78) is illegible to OCR; the sequence proves
%        the page.
%
%  STRUCTURE (from the Indhold, fol. XIV, cross-checked against the body):
%        Forord (fol. III) — Indledning (V–XIII) — Indhold (XIV)
%        A. Viden    pp. 1–22   (a) Videns Begreb 1; b) …; c) Den uendelige
%                                Viden 14, ending in Videns Idee 19)
%        B. Magt     pp. 23–48  (a) Magtens Væsen 23; b) Relativ Magt 30;
%                                c) Den absolute Magt 37, Magtens Idee 41)
%        C. Sandhed  pp. 49–89  (a) Subjectiv Sandhed 50; b) Objectiv
%                                Sandhed 62; c) Den absolute Sandhed 77,
%                                ending in Den logiske Theisme 85)
%
%  CONVENTIONS: diplomatic transcription; 1870 orthography as printed;
%  printer's errors transcribed as printed and logged in % comments at the
%  site; „…“ as printed; hierarchy A./B./C. = \chapter*, a) b) c) =
%  \section*, α) β) γ) = \subsection*. The print is FRAKTUR (long-s visible;
%  Fraktur-tesseract reads it at ~0.8 agreement with the PDF's own text
%  layer, the Latin model at ~0.15) — OCR-first pipeline as in
%  Evangelietroen og Theologien, the GB text layer as second witness.
%
%  STATUS: front matter (Forord fol. III + Indledning V–XIII) and body
%  pp. 1–13 transcribed, image-verified, spliced; pp. 14–89 and the Indhold
%  (fol. XIII–XIV) outstanding. Printer's errors logged so far: meu (VIII),
%  tænkcs (XII), Refleeteren (p.4), Opfatiet (p.6), Idenditets- (p.9),
%  Metpahysiken (p.10) — all transcribed as printed.
% ============================================================
\documentclass[12pt]{book}
\usepackage[utf8]{inputenc}
\usepackage[T1]{fontenc}
\usepackage[danish]{babel}
\usepackage[protrusion=true,expansion=true]{microtype}
\usepackage[margin=1.4in]{geometry}
\usepackage{setspace}
\usepackage{libertinus}
\usepackage{libertinust1math}
\usepackage{textalpha}   % the α) β) γ) subsection marks
\usepackage{fancyhdr}
\usepackage[colorlinks=true,linkcolor=black,urlcolor=black,
  pdftitle={Grundideernes Logik i kort Begreb},
  pdfauthor={Rasmus Nielsen}]{hyperref}
\setlength{\headheight}{15pt}
\pagestyle{fancy}
\fancyhf{}
\fancyhead[LE,RO]{\thepage}
\fancyhead[RE]{\textit{Grundideernes Logik i kort Begreb}}
\fancyhead[LO]{\textit{\leftmark}}
\setlength{\parindent}{1.5em}

\begin{document}

% ---- title page, transcribed from the title leaf proper (PDF 7; the
% ---- ornamental-bordered wrapper leaf, PDF 3 = PDF 5, carries the same
% ---- imprint without the year) ----
\begin{titlepage}
\centering
{\Huge Grundideernes Logik\par}
\vspace{1ex}
{\Large i kort Begreb.\par}
\vspace{4ex}
{\large Af\par}
\vspace{2ex}
{\LARGE R. Nielsen.\par}
\vfill
{\large Kjøbenhavn.\par}
{\large Forlagt af J. H. Schubothes Boghandel.\par}
{\large Louis Kleins Bogtrykkeri.\par}
{\large 1870.\par}
\end{titlepage}

\frontmatter

\chapter*{Forord}
\markboth{Forord}{Forord}
% --- fol. III ---
% (Forord; heading supplied by the skeleton's \chapter*{Forord}.)
I denne lille Bog har Forfatteren bestræbt sig for at give
en kort og klar Fremstilling af sin theistiske Logik. Han har
herved fundet Anledning til at udhæve de Knudepunkter, om
hvilke den videnskabelige Strid mellem Pantheisme og Theisme
væsenlig maa dreie sig.

% signature set in fat-face Fraktur, ranged right
\begin{flushright}R. Nielsen.\end{flushright}

% (fol. IV, PDF 10, is blank.)

\chapter*{Indledning}
\markboth{Indledning}{Indledning}
% print: the heading reads "Indledning." with a short rule beneath.

% --- fol. V ---
% (folio numeral suppressed on this opening page)
Paa Philosophiens nuværende Udviklingstrin bliver det en videnskabelig
Fornødenhed, at behandle Logiken fra tre væsenlig forskjellige Synspunkter.

For at udvikle Tænkningens Love, maa man holde sig til Tænkningen
selv. Men Tænkning er en menneskelig Aandsvirksomhed; vi
maae lære hvad Tænkning er af vor egen Tænken. Det ligger da
nær at betragte Mennesket som det tænkende Subject, besinde sig paa
Tankehandlingernes indbyrdes Forhold og mærke sig de for Forstanden
nødvendige Former, de intellectuelle Spor, som den fremadskridende
Tænkning afsætter. Dette Synspunkt er betegnende for \emph{den
formelle Logik}.

Det er imidlertid saa langt fra, at den endelige Forstand skulde
raade over Tankens Fordringer, at disse tvertimod maae siges at raade
over Tænkeren; der skal tænkes, men naar der tænkes conseqvent, er
det for Resten ligegyldigt, hvem det er, der tænker, om det er et Menneske
eller en Gud eller maaskee den rene Væren. Eftertrykket falder
paa Tænkningen. For Tænkningens Skyld er det derimod af Vigtighed
at fæste Blikket paa det, som tænkes, paa det almeennødvendige
Indhold; thi Tankens Indhold og Tankens Form bestemmes ved hinanden.
Dette Synspunkt er charakteristisk for den \emph{speculative} eller,
for at bruge et maaskee nok saa betegnende Udtryk, \emph{den immanente
Logik}.

Disse Logikens to Hovedformer, af hvilke den første har en Aristoteles,
den sidste en Hegel til øverste Repræsentant, kunne med Grund
siges at nyde en vel fortjent Anerkjendelse i Videnskaben.

At de ikke vel kunne forliges med hinanden, men befinde sig i en
høist betænkelig Strid angaaende Tænkningens allerførste Grundsætninger,
navnlig Modsigelsens Grundsætning, er jo vistnok mindre an%
% --- fol. VI ---
befalende; men de ville begge hævde deres Plads: den ene vil ikke være
istand til at fortrænge den anden.

For at enhver af de stridende Parter kan komme til sin fulde Ret,
er det efter vort Skjøn nødvendigt, at der indføres endnu en tredie
Hovedform. Men førend vi forsøge paa at bevise Rigtigheden af denne
Paastand, ansee vi det for hensigtssvarende at bringe den architektoniske
Modsætning mellem de to Lærebygninger Anskuelsen nærmere ved
at forudskikke et til hver af dem svarende Grundrids.

\section*{Den formelle Logik.}

Den formelle Logik eller Læren om de abstract nødvendige Tankeformer
deler sig i Læren om \emph{Begreb}, \emph{Dom} og \emph{Slutning}.

\subsection*{Begreb.}

Tænkningen, som allerede er med i Sandsningen, hæver sig frem,
idet den opløser de umiddelbare Sandsebilleder, danner ved Sammenligning
og Abstraction et \emph{Almeenbillede}, fastholder det i \emph{Ordet},
klarer sig de nødvendigt sammenhørende Væsensmærker og omdanner det
til et \emph{Begreb} (Begrebsanalyse, Definition).

Indbegrebet af de til et Begreb hørende Bestemmelser udgjør Begrebets
\emph{Indhold}, Antallet af de Gjenstande, hvorpaa det finder Anvendelse,
betegner dets \emph{Omfang}. (Omvendt Forhold mellem Indhold
og Omfang; over- og underordnede Begreber; Slægts- og Artsbegreber).

\subsection*{Dom.}

Begrebet finder kun sit rene logiske Udtryk i Sproget: Ord forbindes
til Sætninger, og i Sætninger udtales Domme. Den Sætning,
hvis Form udtrykker et Begreb, er en \emph{logisk Dom}. For at udtrykke
Begrebet i Formen, maa Dommen formelig angive: \emph{Subjectet}, den
Forestilling, der skal bestemmes, \emph{Prædicatet}, den bestemmende Forestilling,
og \emph{Copula}, det bestemmende Forbindelsesord.

Af Subjectets ved Copula bestemte Forbindelse med Prædicatet
følger, at Dommen kan adskilles i en Flerhed af logisk forskjellige
Domme, og Dommene inddeles: efter \emph{Qvalitet} i bekræftende, benægtende
og uendelige; efter \emph{Qvantitet} i singulære, particulære og generelle;
efter \emph{Relation} i kategoriske, hypothetiske og disjunctive; efter
\emph{Modalitet} i assertoriske, problematiske og apodiktiske.
% --- fol. VII ---

Fremgaaer Prædicatet med almeenlogisk Nødvendighed af Subjectet,
saa er Dommen \emph{analytisk}; kan Prædicatet kun ved en særegen
Anskuelse eller Erfaring føies til Subjectet, saa er Dommen \emph{synthetisk}
(synthetiske Domme \textit{a priori}; Kants Problem).

\subsection*{Slutning.}

Dommen maa begrundes, dens Gyldighed maa indsees; den Tænkningsact,
hvorved en Doms Gyldighed udledes af een eller flere andre
Domme, er en \emph{logisk Slutning}. Har Slutningen kun een Forsætning
(Præmis), kaldes den \emph{umiddelbar}, har den flere Præmisser,
kaldes den \emph{middelbar} (Slutningens Reqvisiter).

De middelbare Slutninger inddeles i: de \emph{kategoriske}, der betinges
af, at det underordnede Begrebs Kjendemærke (\textit{nota notæ}) subsumeres
under det overordnede (syllogistiske Figurer); de \emph{hypothetiske},
der beroe paa det Betingedes Afhængighed af sin Betingelse
(\textit{modus ponens} og \textit{modus tollens}); og de \emph{disjunctive}, som grunde
sig paa, at Leddene af et Hele gjensidig forholde sig saaledes, at naar
eet udtrykkelig er sat, saa ere de øvrige udelukkede (\textit{modus ponendo tollens}),
og at, naar de øvrige ere udelukkede, saa maa det tilbageblevne
sættes (\textit{modus tollendo ponens}).
% The Latin terms (a priori, nota notæ, the modus phrases) are set in
% antiqua against the Fraktur; the Danish "og" between the modus phrases
% is left in Fraktur (upright) here.

Ved at afhandle Læren om Beviis og Methode løser den formelle
Logik sig op i en almindelig Erkjendelses- og Videnskabslære.

\section*{Den immanente Logik.}

Den immanente Logik, der er en dialektisk Udvikling af Kategoriernes
(de al Væren iboende Tankeformers) System, deler sig i Læren
om \emph{Væren}, \emph{Væsen} og \emph{Begreb}.

\subsection*{Væren.}

Den almene Væren er Eet med den almene Tænken; men Tænkens
og Værens Eenhed er i sin allerabstracteste Form indholdsløs: den
rene \emph{Væren} er det Samme som det rene \emph{Intet}. Men Intet er som
Ikke-Væren dog forskjelligt fra Væren. At Væren baade er og ikke er
Intet, vil sige, at Væren uophørlig gaaer over i Intet, og Intet i
Væren: denne Overgang er \emph{Vorden}.

I Vorden er Intet en Grændse for Væren, og Væren for Intet.
Den ved sit Intet begrændsede Væren er et \emph{Værende}, et \emph{Noget}, og
% --- fol. VIII ---
det ved Væren bestemte Intet, et \emph{Andet}. Men Andet er selv Noget,
og Noget igjen Andet; Noget gaaer ved sin egen Grændse over i Andet,
og Andet i Noget: denne Overgang er \emph{Forandring}.

I Forandringen er Noget Eet med sit Andet; dette er det i Eetet
\emph{Uforanderlige}; men Noget er tillige forskjelligt fra Andet; dette er
det i Eetet \emph{Foranderlige}. Det er altsaa det Uforanderlige, der
forandres; meu just derved, at det forandres, holder det sig uforandret:
% printer: »meu« for »men« (so both witnesses; the u confirmed by eye)
under al Forandring er der noget Uforanderligt.

\subsection*{Væsen.}

Det Uforanderlige, der sees i det Foranderlige, er \emph{Væsen}; det
Foranderlige, hvorigjennem Væsenet sees, er \emph{Skin}: Vexelskin af Skin
og Væsen er \emph{Reflexion}.

I sin Identitet med Væsenet er Skinnet \emph{Phænomen}; i sin Forskjel
fra Phænomenet er Væsenet \emph{Grund}. Phænomenet gaaer til
Grunde; det tilhører Skinnet og maa reflecteres i sit Væsen. Men
idet Phænomenet faaer Væsenet til Grund, bliver det begrundet og sat
som \emph{Existens}.

Den i sig reflecterede Existens er et \emph{Existerende}, en \emph{Ting}; men
Tingen kan kun reflecteres i sig ved at reflectere sin Existens i andre
existerende Ting: de til disse Reflexioner svarende Bestemmelser ere
\emph{Egenskaber}.

Under den vexelsidige Reflecteren opløse Tingene sig i Egenskaberne,
og Alt bliver ubestandigt. Men det Ubestandige forudsætter et Bestaaende,
der ikke kan opløses, fordi dets Bestaaen uendelig bekræfter sig
i det Ubestandiges Vexlen; dette Bestaaende er \emph{Substansen}, og de i
Ubestandighed vexlende Fremtrædelsesmaader ere \emph{Accidenserne}.

Forholdet mellem Substansen og dens Accidenser er et \emph{Causalitetsforhold}:
Substansen er \emph{Aarsag}, den oprindelige Sag; Accidenserne
ere \emph{Virkninger}. Men Aarsagen har sit substantielle Udtryk
i Virkningen; uden Virkning ingen Aarsag: ved at gjøre Aarsagen til
Aarsag er Virkningen selv Substans, altsaa en Aarsag, der gjør Aarsagen
til Virkning.

I \emph{Vexelvirkningen} er Virkningen Aarsag, og Aarsagen Virkning;
derved at alle i Væsenet indeholdte Virkelighedsmomenter under
alsidig Vexelvirkning ere udviklede i fuld Totalitet, er det i Totaliteten
udviklede Væsen Begreb.
% --- fol. IX ---

\subsection*{Begreb.}

Den sig begribende Totalitet er i sin Eenhed med sig selv det \emph{Almene},
i sin Forskjel fra sig selv \emph{det Særskilte}, i sin totaliserede
Bestemmethed ved sig selv det \emph{Enkelte}.
% letterspacing as printed: on "det" only before Særskilte, not before
% Almene or Enkelte (verified by zoom)

Begrebets Selvudvikling er dets \emph{Grunddeling} (Dommen). Formedelst
Grunddelingen er Totaliteten i sin første umiddelbare Væren,
det Enkelte, sat som \emph{Subjectet}, og Totaliteten i den reflecterede
Form, det Almene, som \emph{Prædicatet}: Subjectet er, hvad Prædicatet
er; det \emph{Enkelte} er det \emph{Almene}.

Men derved, at det Enkelte begribes som det Almene, og det Almene
som det Enkelte, blive Begrebsadskillelsens modsatte Led gjensidig
begrebne i indbyrdes Vexelbestemmelse som det \emph{Særskilte}.

Grunddelingens Yderled, \emph{Extremerne}, have altsaa under Begrebets
Selvudvikling afsat et \emph{Mellemled}, et \emph{Medium}, hvori de
kunne medieres. Denne Mediation er \emph{Slutningen}.

Begrebets Selvmediation gjennem Dom og Slutning er dets \emph{Subjectivitet};
den ved Mediationen gjenvundne Umiddelbarhed er dets
\emph{Objectivitet}.

Som umiddelbar, tilværende Totalitet er Begrebet Object, og,
ifølge Grunddelingen, en uendelig Mangfoldighed af tilværende Objecter.

Objecterne i deres udvortes, indifferente Gjensidighed ere \emph{mechaniske},
virkende paa hverandre ved \emph{Tryk} og \emph{Stød}. Den herved indtrædende
Adskillelse mellem Bevægelse og Hvile er en relativ Tilsyneladelse
af den \emph{absolute Mechanik}, hvori Bevægelsen er Hvile, og
Hvilen Bevægelse.

Under de mechaniske Objecters Indifferens skjuler sig en \emph{Differens},
der kommer tilsyne i den chemiske Tiltrækning. Objecternes
Differens medfører en Spænding, der udlignes i \emph{det chemiske Product}.
Men Productet, hvori Differensen er neutraliseret, maa som det
\emph{Neutrale} igjen opløse sig i Differens, og \emph{Processen} begynder forfra.

I den chemiske Proces gaaer Begrebet frem af sine egne Forudsætninger,
ophæver Forudsætningerne i deres Resultat, for igjen at ophæve
Resultatet i Forudsætningerne. Det objective Begreb er nu Herre
over Stoffet: det i Forudsætningerne anticiperede Resultat er \emph{Formaal},
Forudsætningerne ere \emph{Midler}, Bevægelsen er \emph{teleologisk}.

Det teleologiske, sig selv i sine Elementer begribende Begreb er en
Form, der giver sig selv Indhold, et Indhold, der giver sig selv Form:
i denne sin \emph{absolute Eenhed} af Form og Indhold, af \emph{Subjectivt
og Objectivt}, af Formaal og Midler, er Begrebet \emph{Idee}.
% --- fol. X ---

Det vil af de foreliggende Grundrids kunne sees, hvor afvigende
Retningerne ere. I den formelle Logik er Mennesket, det tænkende Subject,
fra Først til Sidst beskjæftiget med sine egne Tankehandlinger; det
er Mennesket selv, der abstraherer, reflecterer, dømmer og slutter. Men
just af den Grund ere Tankeformerne tomme Former, og Tankehandlingerne
endelige, psychologisk betingede Forstandsoperationer. Tænkningens
Gjenstand falder udenfor Tænkningen. Her er en uoverstigelig
Kløft mellem det Subjective og det Objective, og den Tænkende har
kun at vaage over, at han i sin Dom om Tingene ikke kommer i Modsigelse
med sig selv. For at sikkre sig Eenheden af det Tænkelige og
det Mulige fastholder man \emph{Modsigelsens Grundsætning}; men
at Modsigelsens Grundsætning ikke afgiver nogen paalidelig Garantie
for den subjective Tænknings objective Gyldighed, have \emph{Antinomierne}
i Kants Kritik af den rene Fornuft tilstrækkelig godtgjort.

Kløften mellem det Subjective og det Objective forsvinder i den
immanente Logik, der abstraherer fra den tænkende Subjectivitet og
lader Væren tænke sig selv ved Hjælp af Modsigelsen og Negationen.
At Væren er Intet, er vel en Modsigelse, thi Væren maa nødvendigviis
udelukke Ikke-Væren. Men den Ikke-Væren, der udelukkes, er netop
den Negation, som ligger i Værens eget Begreb. At det Uforanderlige
er foranderligt, at Væsenet er Skin, at Aarsagen er Virkning, ere lutter
Modsigelser. Men disse Modsigelser gjøre det netop indlysende, hvorledes
den i endelige Forestillinger hildede Tænkning bedrager sig selv
ved at holde det Værendes Indhold udenfor sig. Det Foranderlige er
det Uforanderliges iboende Negation, ligesom Skinnet er Væsenets, og
Virkningen Aarsagens.

Imidlertid tør det dog ikke oversees, at Sagudviklingens immanente
Methode trods sin epochegjørende Betydning dog ogsaa har sine
svage Sider. Negationen er oprindelig en Tankebestemmelse, ikke en
Værensbestemmelse. Det Værende kan i sin Umiddelbarhed hverken negere
sig selv eller et andet Værende. Negationer, Modsigelser og Løsning
af Modsigelser ere Tankehandlinger, der kun gjælde for det Værende
under den udtrykkelige Betingelse, at Væren tænkes. Men for at
denne Betingelse kan opfyldes, maa den paaberaabte Eenhed af Tænken
og Væren fra Begyndelsen af være begrundet paa en tilsvarende Modsætning.
Den rene Væren er Intet, naar Væren tænkes; men Intet
og Væren falde desuagtet ind, hver under sin Sphære: Intet under
Tænkningen, altsaa under det Subjective, og Væren under Virkeligheden,
altsaa under det Objective.

Skal Striden mellem de to anerkjendte Synspunkter ikke blive en%
% --- fol. XI ---
deløs og ufrugtbar, saa maa der vælges et tredie Synspunkt, hvorfra
det bliver muligt at forlige de stridende Parter. Det maa da indrømmes
den formelle Logik, at Tankerne udgaae fra en tænkende Subjectivitet;
men det maa med det Samme indrømmes den immanente
Logik, at det ikke er en endelig, men uendelig Tænkning, der begriber
Tilværelsen. Den formelle Logik maa have Ret deri, at der er en
charakteristisk Forskjel mellem det Subjective i Tænkningen og det Objective
i Virkeligheden; men Immanenslogiken beholder Ret deri, at
Modsætningen af det Subjective og det Objective beroer paa Identitet
af Tænken og Væren. Det er dette tredie Synspunkt, der er valgt i
\emph{Grundideernes Logik}.

Hvad Tankernes Oprindelse angaaer, er der kun tre Antagelser
mulige: 1) at de fremgaae af en endelig Subjectivitet; 2) at de udvikle
sig af den almindelige Væren; 3) at de udspringe af en uendelig,
absolut Subjectivitet. Hvad Forholdet mellem Tænken og Væren angaaer,
er der ligeledes kun tre Antagelser mulige: 1) at det er en endelig
Modsætning, være sig i Overeensstemmelse eller Splid; 2) at
Eenheden ophæver Modsætningen; 3) at Eenheden begrunder Modsætningen.

Til alsidig Udvikling af Logiken ere de tre opstillede Synspunkter
nødvendige; men flere end disse tre ere, hvad Principerne angaaer,
umulige.

% centered rule in the print, closing the first part of the Indledning
\bigskip
\centerline{\rule{6em}{0.4pt}}
\bigskip

% --- fol. XII ---
\section*{Grundideernes Logik}
% print: display heading without full stop; the following text continues
% the sentence begun by the heading.
\noindent udvikler Tankeformernes ideelle Væsenhed under \emph{Videns Idee}, de til
Virkeligheden svarende Værensformer under \emph{Magtens Idee}, og paaviser
det ideelle og det reelle Indholds dualistiske Eenhed i \emph{Sandhedens
Idee}.

Grundadskillelsen af de tre Ideer er særlig motiveret ved Hensyn til
den immanente Logiks Eensidighed.

For at Tænkningen skal være mulig, maa der være Noget, der
tænkcs, altsaa Noget, som Tænkningen selv sætter, forudsætter og frembringer;
% printer: »tænkcs« for »tænkes« (so both witnesses; the c confirmed by eye)
dette er \emph{det Ideelle}. Men for at Virkeligheden, den virkelige
Verden, kan have Gyldighed, maa der være noget Andet end
Tænkningen, Noget, der ikke tænkes, men hvorover der tænkes, Noget,
hvorved Tænkningen bestemmes, men som den ikke selv frembringer;
dette er \emph{det Reelle}. Forskjellen mellem det Ideelle og det Reelle
viser hen til en afgjørende Forskjel mellem Viden og Magt: det Ideelle
er en Vidensbestemmelse, det Reelle en Magtbestemmelse.

Det er almindelig erkjendt, at alt i Verden maa have en Grund;
heri ligger, at det Tilværende maa beroe paa et Ikke-Tilværende.
Grunden for sig er ikke tilværende; den har kun Tilværen i den Existens,
som den begrunder. Grunden er en \emph{Idealitet}; Tingen, den
virkelige Gjenstand, er en \emph{Realitet}. Idealiteten kan tænkes, og der
kan tænkes paa Realiteten; men Realiteten selv kan ikke tænkes; at
tænke Realiteten er at ophæve den og forvandle den til Idealitet.

Den immanente Logik har fortræffelig viist, hvad Idealitet er,
men den har ganske forglemt, hvad Realitet er.

Naar det Endelige, det tilværende Positive, skal holde sig abstract,
som det Reelle ligeoverfor det Uendelige som det Negative, det Ideelle,
da er det let at indsee, at Idealiteten maa opsluge Realiteten. Ved
nemlig at have det Uendelige til Grændse bliver det Endelige selv uendeligt;
men derved bliver det Uendelige, der fra sin Side faaer det Endelige
til Grændse, jo netop endeligt. Idet de modsatte Qvaliteter
saaledes gjensidig have negeret hinanden, have de ved Negationens Ne%
% --- fol. XIII ---
gation med det Samme affirmeret hinanden og derved begge erholdt
samme Idealitet: det Endeliges Idealitet er dets Ophævelse og Bekræftelse
i det Uendelige; ligesaa er det Uendeliges Idealitet dets Ophævelse
og Bekræftelse i det Endelige. Men Realiteten er forsvunden.

Den immanente Logik har en aldeles falsk Opfattelse af Realiteten.
Af Iver for at begribe Gjenstanden forsømmer den at danne sig et
virkeligt Begreb om den virkelige Gjenstand. Den virkelige Gjenstand
er ikke en Idealitet, men en Realitet, et Magtphænomen, som ingen
Tænken formaaer at frembringe. End ikke den guddommelige Tanke
vil være istand til at tænke en virkelig Steen; thi naar Stenen tænkes,
forvandles den til Idealitet og ophører at være til.

Det er med Hensyn hertil, at Grundideernes Logik har indskærpet
den dualistiske Modsætning af Viden og Magt. Skal Magten ligefrem
gaae op i Viden, saa er det ude med Tilværelsens Realitet; Tingene
blive til Skin og den virkelige Verden løser sig op i et metaphysisk
System af guddommelige Tanker. Skal Viden forsvinde i Magten,
saa gaaer Idealiteten tabt. Realismen synker ned til en raa Materialisme,
der gjør det Ideelle til et Skin af det Reelle, det Aandelige
til et tomt Efterskjær af det Legemlige.

Men naar al Værens ideelle Indhold forklares i \emph{Videns Idee},
og Virkelighedens reelle Indhold gaaer op i \emph{Magtens Idee}; saa
følger det jo af sig selv, at Modsætningens Eenhed maa have sit fuldstændige
Udtryk i Ideernes Idee, den alsidige Idee, der mætter Tænkningen
med Idealitet og hævder Virkelighedens Realitet, d.~e. i \emph{Sandhedens
Idee}.

% NB (page-map correction): the Indhold BEGINS on the lower half of this
% same page, fol. XIII (PDF 19) — heading "Indhold.", rule, then A./Viden./
% a) Videns Begreb down to "γ) Den endelige Videns Opløsning … 12." —
% and continues on fol. XIV (PDF 20). It is left for the Indhold batch,
% which is transcribed last against the finished body.

\mainmatter

\chapter*{A. Viden.}
\markboth{A. Viden.}{A. Viden.}
% --- p. 1 ---
% folio suppressed on the opening page; foot of page carries the printer's
% signature line »R. Nielsen: Grundideernes Logik.« with sheet numeral »1«
\section*{a). Videns Begreb.}

At vide Noget er i ubestemt Almindelighed at blive sig
Noget bevidst. Her forudsættes\footnote{Hvad Bestemmelsen af
Philosophiens første almindelige Forudsætning angaaer, maa man vel
skjelne mellem den propædeutiske Forberedelse og den systematiske
Begyndelse. I den philosophiske Propædeutik bliver Forholdet mellem
den Vidende og Videns Gjenstand belyst fra Erkjendelseslærens,
Videnskabslærens og Ideelærens stedse concretere Synspunkter.
(Jvf. Philosophisk Propædeutik i Grundtræk. Kbhvn. 1857). I Logiken
derimod tages Udgangspunktet i den rene, sig selv forudsættende
Viden. (Jvfr. Grundideernes Logik. I. Kbhvn. 1864. Side 1---11).}
% printed footnote mark: 1)
altsaa: \emph{et vidende Selv, en Videns Gjenstand og Bevidsthedens
Eenhed af begge}. Ved Opfattelsen af denne Eenhed kommer det
fornemmelig an paa Vished, Forvisning og Indsigt. Viden er det
intellectuelle Lys, hvori Selvhedsbestemmelser og Bestemmelser af
Andet enes og klares i fuld Gjennemsigtighed.

\subsection*{α) Umiddelbar Vished.}

Ved Sandsning og Tænkning gaaer Tilværelsens mangfoldige Indhold
over i Bevidstheden. Men da Tænkningen i sin Begyndelse ligesaavel
% --- p. 2 ---
er umiddelbar som Sandsningen, saa er Selvvished og Vished om Andet
oprindelig, skjøndt paa forskjellig Maade, umiddelbar.

Den umiddelbare Vished er grundløs, altsaa ubefæstet, og gaaer
uvilkaarlig over i Uvished. Gjenstandene forandre sig: der opkommer
Tvivl om Gjenstandenes Væsen. Sandserne kunne bedrage; der opkommer
Tvivl om Sandsernes Vidnesbyrd. Tænkningen udvikler Begreber; men de
deelviis dannede Begreber staae i Modsigelse med hverandre: der
opkommer Tvivl om Tænkningens Gyldighed
(Skepticisme\footnote{Om den skeptiske Dialektik jvfr. Phil. Propæd.
Side 152---154; Grundid. Log. II. 212---222.}).
% printed footnote mark: 1)

For at Visheden kan gjennem Uvisheden vindes tilbage, udkræves der

\subsection*{β) Forvisning.}

Paa den umiddelbare Visheds Grundlag blev Selvvisheden uden videre
forudsat, og der udviklede sig i Umiddelbarhedens klare Lys en
mangfoldig Viden af Andet (Gjenstandenes og Forestillingernes
Verden). Men Tvivlen kom og fordunklede Klarheden. Tvivlen, der
opløser Viden og forvirrer Bevidstheden, begynder med at undergrave
Visheden om Gjenstandenes Verden og ender med at angribe
Selvvisheden. Men i sit Angreb paa Selvvisheden opløser Tvivlen sig
selv og gaaer under i fornyet Vished. Selv om den Tvivlende kunde
tvivle om Alt, (\textit{de omnibus dubitandum est}) saa kunde han dog
umulig tvivle om sin egen Tvivlen. Tvivlen er Tænken, og Selvet er
tænkende (\textit{cogito ergo sum}), er i sin Selvtænken forvisset om
sin egen Væren. Men derved bliver den tillige forvisset om Tingenes
Væren.

Af Selvforvisningen følger Forvisningen om de oprindelige
Selvhedsformers (de rene Anskuelsers, Kategoriers og Ideers)
subjective Gyldighed (Kant); men deraf følger ogsaa Forvisningen om
deres objective Gyldighed. Dersom
% --- p. 3 ---
der kunde gives Andethedsformer, Værensbestemmelser ved Tingen i sig
(das Ding an sich), der ifølge deres Natur stode i Strid med de
oprindelige Selvhedsformer, da vilde det ikke blot være umuligt for
os, men i sig selv umuligt, at have en Viden om Tingenes Væsen; thi
der kan ikke optages en eneste Bestemmelse i Viden, uden at det
vidende Selv maa frembringe den. Dette gjælder ikke blot om Tanker,
men ogsaa om Sandseindtryk, ikke blot om Begreber, men ogsaa om
Billeder og Forestillinger.

At dømme efter den umiddelbare Erfaring, skulde man jo rigtignok
mene, at der fra Tingenes Verden tilførtes Bevidstheden utallige
Gjenstandsformer og Egenskabsbestemmelser, der aldeles ikke vedkom
dens eget Væsen. Imidlertid vil det ved nogen Eftertanke let vise
sig, at Intet kan vides uden det, som bestemmes i Viden. Hvorledes
skulde der vel kunne gives en Viden om Farven, dersom det vidende
Væsen ikke selv havde Noget med Farven at gjøre.

At beraabe sig paa Synsorganer og Lysindtryk er i denne Sammenhæng
uden Betydning, thi Grundspørgsmaalet er, hvorledes det organisk
betingede Lysindtryk forvandles til Bevidsthed om Lys, eller
hvorledes en Farvebestemmelse kan blive til en Vidensbestemmelse.
Kunde Lyset udenfra skinne gjennem Organet ind i Sjælen, som det
igjennem et Prisme kan falde paa Væggen i Kammeret, da maatte
Farverne jo kunne farve Bevidstheden, saa den blev baade guul og
blaa.

\subsection*{γ) Indsigt.}

Ved Forvisningen bliver den tabte Klarhed gjenvunden i Form af
Opklaring, Begrundelse og Forklaring og svarer til Indsigt. Den
Videndes Indsigt er dobbelt; den er en Indsigt i Sagen med dens Led
og Relationer, og den er en Indsigt i den tænkende Bevidstheds eget
Væsen. Ligesom Tænkning af Andet væsenlig er Selvtænkning, saaledes
er al grundig Forstaaelse af Andet betinget af Selvforstaaelse. Den
logiske Selvforstaaelse beroer paa den Indsigt, at Reflexionens,
Dommens og Slutningens
% --- p. 4 ---
% foot of p. 3 carries the sheet signature »1*«
Former ere almeengyldige Andethedsformer, fordi de ere oprindelige
og nødvendige Selvhedsformer.

Den vidende Subjectivitet er i sin Selvviden et reflecterende Væsen.
Umiddelbart er Reflecteren en Speilen, og forudsætter en speilende
Virken og et Speilbillede: den rene Viden er i intellectuel Forstand
den rene Speilen. At den subjective Virken er Speilen, gjør det
forklarligt, at det Afspeilede kan blive staaende i Andethedens Form
uden at smelte sammen med den bestemmende Selvhedsform, at
Farvebilleder f. Ex. kunne afspeiles af Bevidstheden og vække
Bevidsthed om Farver uden at farve Bevidstheden. Men saa længe
Speilen kun opfattes som umiddelbar
% printer: »Refleeteren« for »Reflecteren« (both witnesses; image confirms ee)
Refleeteren, vil denne Reflecteren ligesaavel kunne passe paa
Andetheden som paa Selvheden. I Optiken er der mange Speile og mange
Slags Speilinger; men logisk forstaaet er den oprindelige Reflecteren
en ideel Selvvirken, en Tænken. Tingene kunne ikke reflectere, thi de
kunne ikke tænke; derimod kan der ikke blot tænkes over deres Væsen,
men dette Væsen kan reflecteres i sig, naar det reflecteres i
Subjectiviteten: al Reflexion af Andet beroer paa Selvreflexion.

Den vidende Subjectivitet er i sin Selvviden en Dom; heri ligger
Grunden til, at Dommens Form er en væsenlig Andethedsform.
Subjectiviteten er en Dom, d.~v.~s. den er et Subject, hvis
Prædicater bestemmes ved de til dens Domme svarende Functioner.
„Selvet er seende, er hørende, er tænkende“; ved Prædicatet „seende“
bestemmes alle Domme om synlige Ting, ved Prædicatet „hørende“ om
hørlige, ved Prædicatet „tænkende“ om tænkelige. Heraf følger, at der
i enhver mulig Dom, hvad enten Selvet forøvrigt dømmer om Andet eller
om sig selv, maa være et dobbelt Subject, Dommens og den Dommendes,
og ligeledes et dobbelt Prædicat og et dobbelt Copula:
Andethedsformen i Dommen er uendelig reflecteret i Selvhedsformen.

Den vidende Subjectivitet er i sin Selvviden en
% --- p. 5 ---
Slutning; heri ligger Grunden til, at Slutningsformen er en væsenlig
Andethedsform. Slutningen fremkommer, forsaavidt der sees paa det
Typiske, derved, at Dommens Subject og Prædicat efter deres
Totalitetsforskjel danne Extremer, hvis Eenhed ikke umiddelbart kan
angives ved Copula, men maa mægles ved et Mellemled.

Selvhedsdommene udsige nemlig, at uligeartede, til forskjellige
Sphærer hørende Prædicater sættes i eet og samme Subject. Ved at see
er Subjectet den Seende, ved at høre er det den Hørende; den Seende
er altsaa hørende, og den Hørende seende; den Sandsende er tænkende,
og den Tænkende sandsende. Hævet over Vidensfunctionernes
Forskjellighed holder Selvet sig fast som det Enkelte. Og dog er
Selvet i sig et Særskilt, thi Sandsning er noget Andet end Tænkning;
hvad der tænkes, sandses ikke: den Sandsende er altsaa ikke den
Tænkende, eller: det sandsende og det tænkende Selv ere særskilte.
Endvidere er Synet noget Andet end Hørelsen; hvad der sees, kan ikke
høres, og hvad der høres, kan ikke sees: den Seende er altsaa ikke
den Hørende, eller: det seende og det hørende Selv ere særskilte.

Selvet vilde nu ved en saadan Deling være adsplittet, hvis Delingen
ikke var en Selvdeling, der lod sig udligne i det Almene. Men Selvet
er selv denne Udligning.

Selvet deler sig i sine Functioner; men disse Functioner, hvor
uligeartede de end ere, ere dog alle Selvhedsfunctioner, og som
Selvhedsfunctioner Functioner af hverandre indbyrdes.

Dersom der i Sandsningen ikke var en skjult Tænkning, saa kunde det
sandsende Selv Intet skjelne, følgelig ikke sandse: Sandsning er
altsaa en Function af Tænkning. Omvendt, dersom der i Tænkningen ikke
var en underforstaaet Sandsning, saa var der intet
Forestillingselement, saa kunde det tænkende Selv Intet skjelne,
følgelig ikke tænke: Tænkning er altsaa en Function af Sandsning. Der
skjelnes i
% --- p. 6 ---
Sandsningen ved Hjælp af Tænkning, og der skjelnes i Tænkningen ved
Hjælp af Sandsning. At der skjelnes, er det Almene; det, som under
den almene Skjelnen adskilles og skjelnes, er det Særskilte. Den
Skjelnende er i sin Skjelnen lig med sig selv: det sig i Alt og Alt i
sig skjelnende Selv er altsaa det Særskilte, fordi det er det Almene.

% printer: »Opfatiet« for »Opfattet« (both witnesses; image shows the i-dot;
% cf. the correct »Opfattet« two paragraphs below)
Opfatiet umiddelbart for sig, er Selvet det Enkelte. I denne sin
Enkelthed er det et fast Punkt, et Atom, hvorfra Virksomheder udgaae
i alle Retninger. Hertil knytter sig Forestillingen om en værende
Ting, et Sjæle-Reale (Herbart), men en saadan Forestilling gjør
Selvet selvløst.

Opfattet umiddelbart i Særskilthedens Form, er Selvet en Vrimmel af
Selvhedspunkter: det er Midtpunktet i den hele Kreds, men ethvert
Punkt i Kredsen er igjen Midtpunkt. Denne uhyre Modsigelse stammer
ikke fra Særskiltheden, men fra den Umiddelbarhed, hvori det
Særskilte holdes fast. Umiddelbarhedens Opløsning i uendelig
Reflexion er Modsigelsens Løsning. Selvet er i sin Enkelthed
reflecteret, thi dets Virken som et Enkelt er formedelst det
Særskilte en Selvvirken. Det er i sin Særskilthed reflecteret, thi
dets særskilte Virken er formedelst dets Enkelthed en Selvvirken:
Selvvirken $=$ Selvvirken er det Almene, hvori Selvet bestaaer.

Viden er ikke Virken, men Lyset i den almene Virken, den uendelige
Reflecteren og Speilen. Den almene Virken, en Virken af Virken i det
Uendelige, er i sin formelle Almindelighed et tomt Begreb. Det i
absolut Forstand almene Væsen er altsaa kun Selvvæsen derved, at det
Enkelte og det Særskilte, hvert for sig i indbyrdes Modsætning, er
Selvvæsen.

Hvor er saa Selvet i dette Væsen? Det er i sin egen Slutning. Naar
det Enkelte og det Særskilte sættes som Extremer, saa er Midten det
Almene (E---A---S); naar det Enkelte og det Almene sættes som
Extremer, saa er Midten
% --- p. 7 ---
det Særskilte (E---S---A); naar endelig det Særskilte og Almene
sættes som Extremer, saa er Midten det Enkelte (A---E---S). Selvet
har Bestaaen i et System af Slutninger: Slutningseenheden, den
absolute Centraleenhed, er selv den evige Begyndelseseenhed.

Fra dette Synspunkt vil det nu være muligt at komme til Indsigt i
Forholdet mellem Selvhedens og Andethedens logiske Grundformer.

Det er i Videns Væsen de samme Former; Modsætningen beroer alene paa
Reflexionens Stilling til det Umiddelbare. I Selvhedsformerne er
Umiddelbarheden fuldstændig opløst; men i Andethedens Former vender
det Umiddelbare idelig tilbage og afbryder Reflexionen. Andetheden
kan ikke reflectere sig selv; den maa reflecteres af Selvheden.

Som Gjenstand for Selvet er det Værende selv et Noget for sig. At det
er Noget for sig selv, tyder paa Selvhed; og dog er Noget kun for sig
selv, hvad det er i sit Forhold til Andet. Men Noget gaaer kun over i
Andet, forsaavidt det tænkes sammen med Andet; dets Overgang er ingen
Selvbevægelse, Bevægelsen er i den tænkende Subjectivitet. Noget
reflecterer sig kun i Andet derved, at Selvet reflecterer paa Noget
og Andet: i Vorden og Forandring viser Andethedsreflexionen sig
udtrykkelig som Selvreflexionens intellectuelle Product og
Modbillede.

Andethedsværen er Vorden. Reflecteres det Værende i sin egen Vorden,
og omvendt: saa bliver Umiddelbarheden ophævet, saa er det Værende
Væsen, og det Vordende Skin. Idet Væsen og Skin reflectere sig i
hinanden, seer det ud, som om Væsenet selv tænkte sit Skin, og
Skinnet sit Væsen; men Tankebevægelsen foregaaer i den Tænkende.
Reflexionen af Skin og Væsen er Selvreflexionens egen Reflexion af
Andet. At den i sig væsenlige Reflexionsform er en Andethedsform,
sees da ogsaa deraf, at Væsenets Reflexion i sig sættes som Ting, og
dets Reflexion i Andet
% --- p. 8 ---
som Egenskab: Tingen selv er selvløs, og Egenskaben en Beskaffenhed
ved det Selvløse.

Men det Selvløse maa nødvendig faae Skin af Selvhed; i modsat Fald
kunde Andethedsformerne umulig blive Vidensformer. Saa veed da det
vidende Selv, at Begrebet af Tingens Væren er en Dom. Tingen bliver
nu Subject, og Egenskaben dens Prædicat. Subjectet er sit Prædicat:
det seer ud i Andetheden ganske som i Selvheden. Men i Andethedens
Væsen er Selvheden Skin; Tingen, det selvløse Subject, har ingen sand
Subjectivitet; Subjectiviteten ligger paa Overfladen, thi den ligger
i Egenskaberne.

Ved at gjøre Tingenes Verden til en Verden af Domme veed det dømmende
Selv at holde Tingene sammen og begribe Sammenhængseenheden i et
System af Slutninger. Sammenhængseenhed bestemmes ved Relation,
Relation ved Reflexion, Reflexion ved Dom, Dom ved Slutning: de
articulerede Vexelbestemmelser af Selvhedens og Andethedens Former
afgive Grundmomenterne i Videns Begreb.

\section*{b) Endelig Viden.}

Ifølge Videns Begreb maa Selvviden og Viden af Andet nødvendig
forudsætte hinanden i det Uendelige; men under den særlige Udvikling
skilles de ad, og heri sees Endeligheden. Særkjendet er, at det
Vidensindhold, der mætter Selvhedsformerne og er nødvendigt til deres
logiske Udfoldning, er fattigt, tomt og formelt i Sammenligning med
det Indhold, der udfylder Andethedens Former med deres udvortes
uendelige Rigdom. Den strenge, systematiske Deduction af de
oprindelige Jeghedssætninger, saaledes som den er bleven gjennemført
i den subjective Idealisme (Fichte), er et talende Exempel paa den
rene Selvvidens aprioriske Fattigdom. I Modsætning hertil har
Erfaringslæren en saadan Overflødighed af Indholdsbestemmelser og
objective Vidensformer, at
% --- p. 9 ---
% a faint speck follows »paavise« (both OCR witnesses read a point);
% judged an ink spot, not a set mark — not transcribed
der ad denne Vei ikke kan være Tale om at paavise den Nødvendighed,
hvormed den objective Verdens Kjendsgjerninger Led for Led fremgaae
af Jeg-Væsenets subjective Grundfordringer (Naturphilosophisk
Mysticisme; Schelling).

For at tydeliggjøre Endelighedsforholdet mærke vi os særlig: først
den aprioriske, dernæst den empiriske Videns Eensidighed og komme
derved til at indsee Nødvendigheden af den endelige Videns logiske
Opløsning.

\subsection*{α) Apriorisk Viden.}

% printer: »Idenditets-« for »Identitets-« (both witnesses; image confirms)
Den formelle Logiks Opfattelse af Idenditets- og
Contradictionsprincipet finder sin adæquateste Anvendelse i den
Videnskab, der netop egner sig til at forstaaeliggjøre den aprioriske
Videns Endelighed, nemlig Mathematiken. Ligningen: A $=$ A, er i
objectiv Forstand, hvad Selvsætningen: Jeg $=$ Jeg, er i subjectiv;
ligeledes har Sætningen: A er ikke Ikke-A, i Andethedens Form,
Sætningen: Jeg er ikke Ikke-Jeg i Selvbegrændsningens Form til
Forudsætning. Her er immanent Reflexion af det Subjective og det
Objective.

Videre kan en hypothetisk Sætning som denne, dersom
$\frac{a}{b} = \frac{c}{x}$, saa er $x = \frac{bc}{a}$, være et
Exempel paa, at objectiv Relation er begrundet i subjectiv Reflexion.
Af Forsætningen maa nemlig Eftersætningen nødvendig følge, Relationen
mellem a, b, c og x er forsaavidt objectiv. Men Størrelserne have
ikke selv sat sig i Relation til hverandre; a har ikke selv divideret
sig med b, eller c med x. Relationen er sat af den reflecterende
Subjectivitet. I Quotienterne $\frac{a}{b}$ og $\frac{c}{x}$ bliver a
reflecteret i b, og b i a, c i x, og x i c; Tælleren forudsætter
Henblik til Nævneren, og Nævneren til Tælleren: det er subjectiv
Reflexion. Vælges en Ligning som: $\frac{dy}{dx} = c$, saa er
Tælleren intet umiddelbart Led, men en Reflex af en i Forsvinden
forsvunden Stør-
% --- p. 10 ---
relse, og det Samme gjælder om Nævneren dx; idet de to Reflexer
reflecteres i hinanden og danne Quotienten, bliver Reflexionen
objectiv; men den er kun objectiv, fordi den er subjectiv.

Kunde den mathematiske Udvikling nu fortsættes saaledes, at der til
enhver ny objectiv Relation svarede et nyt og frugtbart Indblik i den
subjective Reflexion, til enhver ny Læresætning en ny og frugtbar
Erkjendelse af Subjectivitetens oprindelige Selvsætten: da vilde vi
her have en Methode (\textit{methodus mathematica}), som i fortrinlig
Grad egnede sig til at ophæve den aprioriske Videns Endelighed.

Men denne Regning slaaer ikke til. Vel er det sandt, at Sætninger
som, A er A, A er ikke Ikke-A, høre saaledes med i den almindelige
Selvtænkning, at deres Gyldighed er indlysende for enhver Tænkende.
Men hvoraf kommer det, at Ligninger som
$\log A^{x} = x \log A$,
d.\,$\sqrt{a^{2}-x^{2}} = \frac{-\,x\,dx}{\sqrt{a^{2}-x^{2}}}$,
selv om Betegnelserne enkeltviis blive forklarede, ikke kunne siges
at høre med i den oprindelige Begriben af Selvet, ikke ere
uundværlige til logisk Selvforstaaelse? Ja, hvoraf kommer det, at
Mathematiken ikke kan deduceres af
% printer: »Metpahysiken« for »Metaphysiken« (both witnesses; image confirms)
Metpahysiken, eller omvendt? Det kommer deraf, at den aprioriske
Viden, paa Grund af Endeligheden, skiller sig ad i to eensidige
Retninger: en logisk-dialektisk, der fordyber sig i Selvreflexioner,
og en exact-opererende, der taber sig i Reflexioner af
Andet\footnote{Om Forskjellen mellem Reflexions- og
Operationsbegreber jvfr. Mathematik og Dialektik. Kjøbhv. 1859.
S. 40---56.}.
% printed footnote mark: 1)

\subsection*{β) Empirisk Viden.}

De aprioriske Andethedsformer, der articuleres i de rene
Størrelsesanalyser, støtte sig alene til Phantasie-Anskuelser og
Forestillinger, og den forestillede Andethed falder endnu forsaavidt
indenfor Selvheden.

Men naar det gjælder en Viden om existerende Ting,
% --- p. 11 ---
indtræder den Modsigelse, at Gjenstandene indenfor Viden selv falde
udenfor Viden. Den endelige Bevidsthed har nemlig en Verden af
Gjenstande \emph{udenfor} sig; men den kan umulig have en Viden om
disse Gjenstande uden ved at optage deres Væren i sig, og saaledes
have dem \emph{indenfor} sig. Denne i den empiriske Videns Grundlag
skjulte Modsigelse betinger Erkjendelsens Mangfoldighed, men
forraader tillige dens Endelighed under Form af Tilfældighed.

Tilfældigheden er subjectiv, thi de almeen-nødvendige
Andethedsformer, hvori Gjenstandene maae opfattes, kunne jo haves i
Bevidstheden, uden Viden om, at de særlige Gjenstande virkelig
existere. At en tilværende Gjenstand altid maa svare til et eller
andet Begreb, d.~v.~s. til en eller anden Udfyldning af aprioriske
Former, ophæver ikke Tilfældigheden, da en saadan Udfyldning kan skee
paa utallige Maader: at mange og forskjellige Ting bære samme Navn,
tyder netop paa, at Forskjellighederne ansees for tilfældige. Det
empiriske Begreb, hvorpaa Navnet lyder, gjentager sig i en Mængde
tilsvarende Exemplarer; paa Begrebets Grundlag bliver Exemplaret
beskrevet, men Beskrivelsens billedlige Elementer angive
Tilfældigheden.

I Forhold til Individerne og det umiddelbart Individuelle er det
tilværende Begreb Arten, og de særskilte Arters nærmeste Almeenbegreb
er Slægten. At et Individ med et eller flere Mærker af en vis Art maa
have alle til Arten hørende Mærker, sluttes ved \emph{Analogie}; at
Væsensbestemmelser, der ere afgjørende for nogle under en vis Art
eller Slægt hørende Individer, maae gjælde for dem alle, sluttes ved
\emph{Induction}; men da ikke blot Individerne men ogsaa Begreberne
ere empiriske, saa er saavel Inductionen som Analogien behæftet med
et Tilfældighedsmoment, der kan forvandle Visheden til Uvished.

Fra et almindeligt Synspunkt kan den empiriske Viden vel opnaae
tilforladelig Vished; men de særskilte Tilfælde
% --- p. 12 ---
krydse bestandig det Almene og fremkalde Uvished. Forvisningsgrunden
er negativ: „her er intet Tilfælde, der gjør Brud paa det ved
Iagttagelserne fastsatte Almene; ergo er dette Tilfældene bestemmende
Almene en Lov“. Men det saaledes fastsatte Almene er dog i sin
positive Form empirisk. For de methodisk fortsatte Iagttagelser
frembyder Erfaringsverdenen bestandig nye Sider, nye Relationer, nye
Egenskaber, nye Love.

Tilfældigheden er objectiv, thi det er ikke blot os, der maae
forandre, corrigere og forbedre vore Begreber om Tingene, men selv
Begreberne (Arter og Slægter; Darwin) forandre sig efterhaanden, idet
de som Existensbegreber forandre det existentielle Udtryk med
Individernes Forandringer. Tilfældigheden er objectiv, thi
Gjenstandenes Existens er uafhængig af vor subjective Viden om dem.
Vi kunde ikke sandse Gjenstandene, dersom Gjenstandene ikke vare
sandselig til; men de sandselige Gjenstande kunne meget vel være til,
uden at vi sandse dem. Vi danne os Begreber, der abstraheres af
Gjenstandene, men Gjenstandene selv ere ligegyldige mod vor
Abstraheren og Begriben. Det ligger i Gjenstandenes Selvstændighed at
have vor Viden om dem udenfor sig som en Tilfældighed.

For vor endelige Viden staaer altsaa Tingenes Verden som en Skranke,
et uigjennemtrængeligt Andet. Vi erkjende, siger Erfaringslæren,
Tingenes Egenskaber, men ikke Tingene selv; vi erkjende Lovene, men
ikke Væsenet. Mechaniken veed ikke, hvad Kraft er; den chemiske
Physik veed ikke, hvad Elektricitet er; Physiologien veed ikke, hvad
Livet er. Den sande empiriske Viden har overalt i sine Grundlag
optaget Bestemmelsen: Ikke-Viden.

\subsection*{γ) Den endelige Videns Opløsning.}

De Modsigelser, hvormed den endelige Viden fra alle Sider er
behæftet, samle sig tilsidst i den ene: Alvidenhed er nødvendig og
dog ubegribelig.
% --- p. 13 ---
At Alvidenhed er nødvendig, fremlyser af Videns oprindelige Væsen. Da
nemlig alle Tilværelsesbestemmelser ere i logisk Forstand
Vidensbestemmelser, og alle Vidensbestemmelser Functioner af det
vidende Selv: saa kan slet Intet vides, uden alene det, som er
frembragt ved Viden.

Modsigelsen er nu den, at den endelige Viden overalt i den virkelige
Verden forefinder Gjenstandsformer, som hverken den selv eller
Gjenstanden har frembragt. Alle udvortes Ting ere til i Rummet. Men
hvad er Rummet? En ved Anskuen frembragt, for Anskuelsen værende
Udstrækning. Idet vi anskue Tingene i Rummet, anskue vi dem i en
forud anskuet Form (det objective Rum); og det Samme gjælder om
Tiden.

Ved at reflectere paa de i Rummet og Tiden tilværende Gjenstande,
komme vi til et lignende Resultat. For os med vor endelige Viden ere
de anskuelige Gjenstande til, før vi anskue dem. Men Anskueligt er
kun det, som er frembragt ved Anskuen. Vor Anskuen bestemmes ved det
Anskuelige, det allerede Frembragte. Hvorledes er Gjenstanden bleven
anskuelig? Alene derved, at dens Tilværelse er sat i en anskuelig,
d.~e. i en for Anskuelsen mulig Form; men en for Anskuelsen mulig
Gjenstandsform er kun virkelig i Kraft af en oprindelig, Gjenstanden
formende Anskuen.

For den oprindelige Anskuen er Gjenstanden selv reflecteret som
Billede, og Billedet som Gjenstand. I Billedets Form er
Gjenstandsreflexionen imidlertid overfladisk; det væsenlige Indhold
har sin udtømmende Form i Begrebet. For vor Viden er det Gjenstanden,
der bestemmer Begrebet; men i Sagen selv er det Begrebet, der
bestemmer Gjenstanden. For at bestemme Gjenstanden, maa Begrebet
forud være bestemt ved en af vort Begreb uafhængig Begriben, altsaa
ved en af vor Viden uafhængig Viden.

Den Alvidenhedens Anskuen, Tænken og Begriben, hvori Tilværelsen
skulde have sin oprindelige Grund, er jo vistnok
% (sentence continues on p. 14)
% [text to be added: pp. 14--22]

\chapter*{B. Magt.}
\markboth{B. Magt.}{B. Magt.}
% [text to be added: pp. 23--36]
% [text to be added: pp. 37--48]

\chapter*{C. Sandhed.}
\markboth{C. Sandhed.}{C. Sandhed.}
% [text to be added: pp. 49--61]
% [text to be added: pp. 62--76]
% [text to be added: pp. 77--89]

\backmatter

\chapter*{Indhold}
\markboth{Indhold}{Indhold}
% [text to be added: fol. XIV (Indhold, transcribed last against the finished body)]

\end{document}
